\section{Phase 10: A Full Block-Banded Space-Time Optimization Framework Based on Raw Scheme C}
\label{sec:phase10_full_framework}

The purpose of this phase is no longer to determine whether BLOM resembles the global MINCO reference---that is the role of Phases~7--8---nor merely to verify that a minimal differentiable loop exists---that is the role of Phase~9.
Instead, the goal here is to elevate raw Scheme~C into a complete block-banded space-time optimization framework.

The central optimization variable remains
\[
\theta=(\bar q,T),
\qquad
\bar q=(q_1,\dots,q_{M-1})^\top\in\mathbb R^{M-1},
\qquad
T=(T_1,\dots,T_M)^\top\in\mathbb R_{>0}^{M},
\]
and the trajectory coefficients are generated by the raw Scheme~C map
\[
c^{\mathrm{C},k}(\theta)\in\mathbb R^{2sM}.
\]
As in the previous phases, the spline is scalar-valued in the canonical presentation and each segment coefficient block is
\[
c_i(\theta)\in\mathbb R^{2s}.
\]
Although the implementation target of this phase includes general even \(k\), the mathematical statements below will make explicit whenever a result is conditional on general-\(k\) properties that are not yet fully proved.

\subsection{Admissible domain and time reparameterization}
\label{subsec:phase10_time_reparam}

To enforce strict positivity of the segment durations during optimization, we introduce an unconstrained time variable
\[
\tau=(\tau_1,\dots,\tau_M)^\top\in\mathbb R^M
\]
and define
\begin{equation}
\label{eq:phase10_softplus_time}
T_i(\tau_i):=T_{\min}+\operatorname{softplus}(\tau_i),
\qquad
\operatorname{softplus}(x):=\log(1+e^x),
\qquad
T_{\min}>0.
\end{equation}
The new optimization variable is
\[
\xi:=(\bar q,\tau)\in\mathbb R^{M-1}\times\mathbb R^M.
\]
The induced duration map is written compactly as
\[
T(\tau)=\bigl(T_1(\tau_1),\dots,T_M(\tau_M)\bigr)^\top.
\]
We shall optimize the reparameterized objective
\[
\widetilde J(\xi):=J(\bar q,T(\tau)).
\]

\begin{proposition}[Smooth positivity-preserving time parameterization]
\label{prop:phase10_softplus}
The map \(\tau\mapsto T(\tau)\) defined by \eqref{eq:phase10_softplus_time} is \(C^\infty\), satisfies
\[
T_i(\tau_i)>T_{\min}>0
\qquad
\forall \tau_i\in\mathbb R,
\]
and has derivative
\begin{equation}
\label{eq:phase10_softplus_derivative}
\frac{dT_i}{d\tau_i}
=
\sigma(\tau_i)
:=
\frac{1}{1+e^{-\tau_i}}
\in(0,1).
\end{equation}
\end{proposition}

\begin{proof}
The softplus function is \(C^\infty\) on \(\mathbb R\), strictly positive, and satisfies
\[
\frac{d}{dx}\operatorname{softplus}(x)=\frac{e^x}{1+e^x}=\sigma(x).
\]
Therefore each component \(T_i(\tau_i)\) is \(C^\infty\), strictly greater than \(T_{\min}\), and has derivative \eqref{eq:phase10_softplus_derivative}.
\end{proof}

\subsection{A full objective layer}
\label{subsec:phase10_full_objective}

The optimization objective of Phase~10 is written as
\begin{equation}
\label{eq:phase10_full_K}
K(c,T)
=
K_{\mathrm{ctrl}}(c,T)
+
\lambda_T K_{\mathrm{time}}(T)
+
\lambda_{\mathrm{obs}} K_{\mathrm{obs}}(c,T)
+
\lambda_{\mathrm{dyn}} K_{\mathrm{dyn}}(c,T)
+
\lambda_{\mathrm{bc}} K_{\mathrm{bc}}(c,T)
+
\lambda_{\mathrm{reg}} K_{\mathrm{reg}}(T),
\end{equation}
with nonnegative scalar weights
\[
\lambda_T,\lambda_{\mathrm{obs}},\lambda_{\mathrm{dyn}},\lambda_{\mathrm{bc}},\lambda_{\mathrm{reg}}\ge 0.
\]
The total reduced objective is
\begin{equation}
\label{eq:phase10_J}
J(\bar q,T):=K(c^{\mathrm{C},k}(\bar q,T),T),
\qquad
\widetilde J(\bar q,\tau):=J(\bar q,T(\tau)).
\end{equation}

To define the sampled penalties rigorously, recall the segment polynomial basis from the previous phases.
For each derivative order \(r\ge 0\), let
\[
B^{(r)}(\tau)\in\mathbb R^{1\times 2s}
\]
denote the row operator mapping a coefficient block \(c_i\) to the \(r\)-th derivative evaluated at local time \(\tau\):
\[
p_i^{(r)}(\tau)=B^{(r)}(\tau)c_i.
\]
For a finite derivative index set
\[
\mathcal R=\{r_1,\dots,r_m\},
\]
define the stacked sampling operator
\begin{equation}
\label{eq:phase10_stacked_B}
\mathcal B_{\mathcal R}(\tau)
:=
\begin{bmatrix}
B^{(r_1)}(\tau)\\
\vdots\\
B^{(r_m)}(\tau)
\end{bmatrix}
\in\mathbb R^{m\times 2s}.
\end{equation}
Then the sampled local jet is
\begin{equation}
\label{eq:phase10_sampled_jet}
z_{i,\nu}^{\mathcal R}(c_i,T_i)
:=
\mathcal B_{\mathcal R}(\alpha_\nu T_i)c_i,
\qquad
0\le \alpha_\nu\le 1.
\end{equation}

\subsubsection{Control effort}
\label{subsubsec:phase10_ctrl}

As in Phase~1, let \(H_i(T_i)\in\mathbb R^{2s\times 2s}\) be the segment Gram matrix of the \(s\)-th derivative control cost.
Define
\begin{equation}
\label{eq:phase10_Kctrl}
K_{\mathrm{ctrl}}(c,T)
:=
\frac12\sum_{i=1}^{M} c_i^\top H_i(T_i)c_i.
\end{equation}

\subsubsection{Time regularization}
\label{subsubsec:phase10_time}

Let \(\psi_T:(T_{\min},\infty)\to\mathbb R\) be \(C^1\).
Define
\begin{equation}
\label{eq:phase10_Ktime}
K_{\mathrm{time}}(T)
:=
\sum_{i=1}^{M}\psi_T(T_i).
\end{equation}
The minimal choice used in Phase~9 is recovered by \(\psi_T(T)=T\).

\subsubsection{Obstacle penalty}
\label{subsubsec:phase10_obs}

Fix a finite set of sample locations
\[
0\le \alpha_1^{\mathrm{obs}},\dots,\alpha_{N_{\mathrm{obs}}}^{\mathrm{obs}}\le 1,
\qquad
w_\nu^{\mathrm{obs}}>0,
\]
and a \(C^1\) obstacle penalty
\[
\Phi_{\mathrm{obs},\nu}:\mathbb R^{m_{\mathrm{obs}}}\times (T_{\min},\infty)\to\mathbb R_{\ge 0},
\]
where \(m_{\mathrm{obs}}\) is determined by the selected derivative index set \(\mathcal R_{\mathrm{obs}}\).
Define
\begin{equation}
\label{eq:phase10_Kobs}
K_{\mathrm{obs}}(c,T)
:=
\sum_{i=1}^{M}\sum_{\nu=1}^{N_{\mathrm{obs}}}
w_\nu^{\mathrm{obs}}\,
\Phi_{\mathrm{obs},\nu}
\bigl(
z_{i,\nu}^{\mathcal R_{\mathrm{obs}}}(c_i,T_i),\,T_i
\bigr).
\end{equation}
The scalar Phase~9 soft obstacle term is the special case \(\mathcal R_{\mathrm{obs}}=\{0\}\).

\subsubsection{Dynamics penalty}
\label{subsubsec:phase10_dyn}

Similarly, choose sample locations
\[
0\le \alpha_1^{\mathrm{dyn}},\dots,\alpha_{N_{\mathrm{dyn}}}^{\mathrm{dyn}}\le 1,
\qquad
w_\nu^{\mathrm{dyn}}>0,
\]
and a \(C^1\) dynamics penalty
\[
\Phi_{\mathrm{dyn},\nu}:\mathbb R^{m_{\mathrm{dyn}}}\times (T_{\min},\infty)\to\mathbb R_{\ge 0}
\]
defined on the sampled jet associated with \(\mathcal R_{\mathrm{dyn}}\).
Define
\begin{equation}
\label{eq:phase10_Kdyn}
K_{\mathrm{dyn}}(c,T)
:=
\sum_{i=1}^{M}\sum_{\nu=1}^{N_{\mathrm{dyn}}}
w_\nu^{\mathrm{dyn}}\,
\Phi_{\mathrm{dyn},\nu}
\bigl(
z_{i,\nu}^{\mathcal R_{\mathrm{dyn}}}(c_i,T_i),\,T_i
\bigr).
\end{equation}

\subsubsection{Boundary penalty}
\label{subsubsec:phase10_bc}

Let
\[
\mathcal R_{\mathrm{bc}}\subset\{0,\dots,s-1\}
\]
be the set of endpoint jet orders to be softly enforced.
Let the desired initial and terminal jets be
\[
d^- \in \mathbb R^{|\mathcal R_{\mathrm{bc}}|},
\qquad
d^+ \in \mathbb R^{|\mathcal R_{\mathrm{bc}}|}.
\]
Let
\[
\Phi_{\mathrm{bc}}^-:\mathbb R^{|\mathcal R_{\mathrm{bc}}|}\to\mathbb R_{\ge 0},
\qquad
\Phi_{\mathrm{bc}}^+:\mathbb R^{|\mathcal R_{\mathrm{bc}}|}\times (T_{\min},\infty)\to\mathbb R_{\ge 0}
\]
be \(C^1\).
Define
\begin{equation}
\label{eq:phase10_Kbc}
K_{\mathrm{bc}}(c,T)
:=
\Phi_{\mathrm{bc}}^-\!\left(
\mathcal B_{\mathcal R_{\mathrm{bc}}}(0)c_1-d^-
\right)
+
\Phi_{\mathrm{bc}}^+\!\left(
\mathcal B_{\mathcal R_{\mathrm{bc}}}(T_M)c_M-d^+,\,T_M
\right).
\end{equation}

\subsubsection{Additional regularization}
\label{subsubsec:phase10_reg}

For concreteness, we use a \(C^1\) time-smoothing regularizer
\begin{equation}
\label{eq:phase10_Kreg}
K_{\mathrm{reg}}(T)
:=
\frac12\sum_{i=1}^{M-1}(T_{i+1}-T_i)^2.
\end{equation}
Other \(C^1\) regularizers may be added without changing the chain-rule structure below.

\subsection{Differentiability of the full objective layer}
\label{subsec:phase10_diff}

\begin{proposition}[Differentiability of the full objective layer]
\label{prop:phase10_K_C1}
Assume:
\begin{enumerate}
    \item each segment Gram map \(T_i\mapsto H_i(T_i)\) is \(C^1\);
    \item \(\psi_T\in C^1\);
    \item all sampled penalties \(\Phi_{\mathrm{obs},\nu}\), \(\Phi_{\mathrm{dyn},\nu}\), \(\Phi_{\mathrm{bc}}^\pm\) are \(C^1\);
    \item \(K_{\mathrm{reg}}\) is \(C^1\).
\end{enumerate}
Then \(K(c,T)\) defined by \eqref{eq:phase10_full_K} is \(C^1\) in \((c,T)\).
\end{proposition}

\begin{proof}
The control term is \(C^1\) because it is a finite sum of functions of the form
\[
(c_i,T_i)\mapsto \frac12 c_i^\top H_i(T_i)c_i
\]
with \(H_i\in C^1\).
The time term is \(C^1\) because \(\psi_T\in C^1\).
For the obstacle and dynamics terms, each sampled jet
\[
(c_i,T_i)\mapsto z_{i,\nu}^{\mathcal R}(c_i,T_i)
=
\mathcal B_{\mathcal R}(\alpha_\nu T_i)c_i
\]
is \(C^1\), since \(\mathcal B_{\mathcal R}\) is polynomial in its scalar argument.
Composing with the \(C^1\) penalties yields \(C^1\) summands, hence \(K_{\mathrm{obs}}\) and \(K_{\mathrm{dyn}}\) are \(C^1\).
The boundary term is treated identically: both \(\mathcal B_{\mathcal R_{\mathrm{bc}}}(0)c_1\) and \(\mathcal B_{\mathcal R_{\mathrm{bc}}}(T_M)c_M\) are \(C^1\), so \(K_{\mathrm{bc}}\in C^1\).
Finally, \(K_{\mathrm{reg}}\in C^1\) by assumption.
Therefore \(K\in C^1\).
\end{proof}

\subsection{Block derivatives of the full objective}
\label{subsec:phase10_block_grads}

Let
\[
g_c(c,T):=\frac{\partial K}{\partial c}(c,T)
=
\begin{bmatrix}
g_{c,1}(c,T)\\
\vdots\\
g_{c,M}(c,T)
\end{bmatrix},
\qquad
g_{c,i}(c,T)\in\mathbb R^{2s},
\]
and
\[
g_T(c,T):=\frac{\partial K}{\partial T}(c,T)
=
\begin{bmatrix}
g_{T,1}(c,T)\\
\vdots\\
g_{T,M}(c,T)
\end{bmatrix}
\in\mathbb R^{M}.
\]

For each derivative index set \(\mathcal R\), define the derivative of the stacked operator
\[
\mathcal B_{\mathcal R}'(\tau)
:=
\frac{d}{d\tau}\mathcal B_{\mathcal R}(\tau).
\]

\begin{proposition}[Exact block formula for \texorpdfstring{$\partial K/\partial c$}{dK/dc}]
\label{prop:phase10_dKdc}
For each segment \(i\),
\begin{align}
\label{eq:phase10_dKdc}
g_{c,i}
&=
H_i(T_i)c_i \notag\\
&\quad
+\lambda_{\mathrm{obs}}
\sum_{\nu=1}^{N_{\mathrm{obs}}}
w_\nu^{\mathrm{obs}}
\,
\mathcal B_{\mathcal R_{\mathrm{obs}}}(\alpha_\nu^{\mathrm{obs}}T_i)^\top
\nabla_1 \Phi_{\mathrm{obs},\nu}
\bigl(
z_{i,\nu}^{\mathcal R_{\mathrm{obs}}},T_i
\bigr)
\notag\\
&\quad
+\lambda_{\mathrm{dyn}}
\sum_{\nu=1}^{N_{\mathrm{dyn}}}
w_\nu^{\mathrm{dyn}}
\,
\mathcal B_{\mathcal R_{\mathrm{dyn}}}(\alpha_\nu^{\mathrm{dyn}}T_i)^\top
\nabla_1 \Phi_{\mathrm{dyn},\nu}
\bigl(
z_{i,\nu}^{\mathcal R_{\mathrm{dyn}}},T_i
\bigr)
\notag\\
&\quad
+\lambda_{\mathrm{bc}}\,\delta_{i1}\,
\mathcal B_{\mathcal R_{\mathrm{bc}}}(0)^\top
\nabla \Phi_{\mathrm{bc}}^-
\!\left(
\mathcal B_{\mathcal R_{\mathrm{bc}}}(0)c_1-d^-
\right)
\notag\\
&\quad
+\lambda_{\mathrm{bc}}\,\delta_{iM}\,
\mathcal B_{\mathcal R_{\mathrm{bc}}}(T_M)^\top
\nabla_1 \Phi_{\mathrm{bc}}^+
\!\left(
\mathcal B_{\mathcal R_{\mathrm{bc}}}(T_M)c_M-d^+,\,T_M
\right).
\end{align}
Here \(\delta_{ij}\) is the Kronecker delta and \(\nabla_1\) denotes the gradient with respect to the first argument.
\end{proposition}

\begin{proof}
Differentiate the control term blockwise:
\[
\frac{\partial}{\partial c_i}
\left(\frac12 c_i^\top H_i(T_i)c_i\right)=H_i(T_i)c_i,
\]
since \(H_i(T_i)\) is symmetric.
For a generic sampled penalty
\[
\Phi(z_{i,\nu}^{\mathcal R}(c_i,T_i),T_i),
\]
the chain rule gives
\[
\frac{\partial}{\partial c_i}\Phi(z_{i,\nu}^{\mathcal R},T_i)
=
\left(\frac{\partial z_{i,\nu}^{\mathcal R}}{\partial c_i}\right)^\top
\nabla_1 \Phi(z_{i,\nu}^{\mathcal R},T_i).
\]
Because
\[
z_{i,\nu}^{\mathcal R}
=
\mathcal B_{\mathcal R}(\alpha_\nu T_i)c_i,
\]
we have
\[
\frac{\partial z_{i,\nu}^{\mathcal R}}{\partial c_i}
=
\mathcal B_{\mathcal R}(\alpha_\nu T_i).
\]
This yields the obstacle and dynamics contributions.
The boundary term contributes only for \(i=1\) and \(i=M\), producing the last two lines of \eqref{eq:phase10_dKdc}.
The time and regularization terms do not depend on \(c_i\).
Summing all contributions proves the formula.
\end{proof}

\begin{proposition}[Exact block formula for \texorpdfstring{$\partial K/\partial T$}{dK/dT}]
\label{prop:phase10_dKdT}
For each \(i=1,\dots,M\),
\begin{align}
\label{eq:phase10_dKdT}
g_{T,i}
&=
\frac12 c_i^\top H_i'(T_i)c_i
+
\lambda_T \psi_T'(T_i)
+
\lambda_{\mathrm{reg}}
\Bigl[(2T_i-T_{i-1}-T_{i+1})\Bigr]_{\mathrm{valid}}
\notag\\
&\quad
+\lambda_{\mathrm{obs}}
\sum_{\nu=1}^{N_{\mathrm{obs}}}
w_\nu^{\mathrm{obs}}
\Bigl[
\partial_2 \Phi_{\mathrm{obs},\nu}
\bigl(
z_{i,\nu}^{\mathcal R_{\mathrm{obs}}},T_i
\bigr)
+
\alpha_\nu^{\mathrm{obs}}
\bigl\langle
\nabla_1 \Phi_{\mathrm{obs},\nu}
\bigl(
z_{i,\nu}^{\mathcal R_{\mathrm{obs}}},T_i
\bigr),
\,
\mathcal B'_{\mathcal R_{\mathrm{obs}}}(\alpha_\nu^{\mathrm{obs}}T_i)c_i
\bigr\rangle
\Bigr]
\notag\\
&\quad
+\lambda_{\mathrm{dyn}}
\sum_{\nu=1}^{N_{\mathrm{dyn}}}
w_\nu^{\mathrm{dyn}}
\Bigl[
\partial_2 \Phi_{\mathrm{dyn},\nu}
\bigl(
z_{i,\nu}^{\mathcal R_{\mathrm{dyn}}},T_i
\bigr)
+
\alpha_\nu^{\mathrm{dyn}}
\bigl\langle
\nabla_1 \Phi_{\mathrm{dyn},\nu}
\bigl(
z_{i,\nu}^{\mathcal R_{\mathrm{dyn}}},T_i
\bigr),
\,
\mathcal B'_{\mathcal R_{\mathrm{dyn}}}(\alpha_\nu^{\mathrm{dyn}}T_i)c_i
\bigr\rangle
\Bigr]
\notag\\
&\quad
+\lambda_{\mathrm{bc}}\,\delta_{iM}
\Bigl[
\partial_2 \Phi_{\mathrm{bc}}^+
\!\left(
\mathcal B_{\mathcal R_{\mathrm{bc}}}(T_M)c_M-d^+,\,T_M
\right)
+
\bigl\langle
\nabla_1 \Phi_{\mathrm{bc}}^+
\!\left(
\mathcal B_{\mathcal R_{\mathrm{bc}}}(T_M)c_M-d^+,\,T_M
\right),
\,
\mathcal B'_{\mathcal R_{\mathrm{bc}}}(T_M)c_M
\bigr\rangle
\Bigr].
\end{align}
Here \([\cdot]_{\mathrm{valid}}\) denotes the standard tridiagonal time-smoothing contribution, with the obvious endpoint simplifications:
\[
\frac{\partial K_{\mathrm{reg}}}{\partial T_1}=T_1-T_2,
\qquad
\frac{\partial K_{\mathrm{reg}}}{\partial T_M}=T_M-T_{M-1},
\]
and for \(2\le i\le M-1\),
\[
\frac{\partial K_{\mathrm{reg}}}{\partial T_i}=2T_i-T_{i-1}-T_{i+1}.
\]
\end{proposition}

\begin{proof}
Differentiate each term with respect to \(T_i\).
For the control term:
\[
\frac{\partial}{\partial T_i}
\left(\frac12 c_i^\top H_i(T_i)c_i\right)
=
\frac12 c_i^\top H_i'(T_i)c_i.
\]
For the time term:
\[
\frac{\partial}{\partial T_i}\lambda_T\sum_{j=1}^{M}\psi_T(T_j)
=
\lambda_T\psi_T'(T_i).
\]
For the regularizer \eqref{eq:phase10_Kreg}, direct differentiation gives the tridiagonal formula above.

For a generic sampled term
\[
\Phi(z_{i,\nu}^{\mathcal R}(c_i,T_i),T_i),
\]
the chain rule yields
\[
\frac{\partial}{\partial T_i}\Phi(z_{i,\nu}^{\mathcal R},T_i)
=
\partial_2 \Phi(z_{i,\nu}^{\mathcal R},T_i)
+
\left\langle
\nabla_1 \Phi(z_{i,\nu}^{\mathcal R},T_i),
\,
\frac{\partial z_{i,\nu}^{\mathcal R}}{\partial T_i}
\right\rangle.
\]
Now
\[
z_{i,\nu}^{\mathcal R}
=
\mathcal B_{\mathcal R}(\alpha_\nu T_i)c_i,
\]
hence
\[
\frac{\partial z_{i,\nu}^{\mathcal R}}{\partial T_i}
=
\alpha_\nu \mathcal B_{\mathcal R}'(\alpha_\nu T_i)c_i.
\]
Substituting this identity yields the obstacle and dynamics contributions.
The initial boundary term does not depend on any \(T_i\), while the terminal boundary term depends only on \(T_M\), yielding the last line of \eqref{eq:phase10_dKdT}.
Combining all contributions proves the result.
\end{proof}

\subsection{General-\texorpdfstring{$k$}{k} raw Scheme C map and block-banded reverse differentiation}
\label{subsec:phase10_general_k_backward}

At this stage, the complete general-\(k\) theory of the raw Scheme~C map is not yet fully proved.
Accordingly, we state the exact assumptions needed to derive the reverse-differentiation framework.

\begin{assumption}[General-\texorpdfstring{$k$}{k} differentiable raw Scheme C map]
\label{ass:phase10_general_k_C1}
For each even \(k\in 2\mathbb N\), there exists an admissible domain
\[
\Theta_{\mathrm{adm}}^{(k)}\subset \mathbb R^{M-1}\times\mathbb R_{>0}^{M}
\]
such that the raw Scheme~C coefficient map
\[
c^{\mathrm{C},k}:\Theta_{\mathrm{adm}}^{(k)}\to \mathbb R^{2sM}
\]
is \(C^1\).
\end{assumption}

\begin{remark}
\label{rem:phase10_general_k_C1_status}
A full proof of Assumption~\ref{ass:phase10_general_k_C1} for all even \(k\) is \textbf{not completed at the present stage}.
Phase~9 rigorously established the canonical \(k=2\) case.
Phase~10 treats general \(k\) as the correct framework-level target.
\end{remark}

\begin{assumption}[General-\texorpdfstring{$k$}{k} block-local support]
\label{ass:phase10_general_k_support}
For each segment \(i\in\{1,\dots,M\}\) and each even \(k\), there exist support sets
\[
\mathcal S_q(i,k)\subset \{1,\dots,M-1\},
\qquad
\mathcal S_T(i,k)\subset \{1,\dots,M\},
\]
such that:
\begin{equation}
\label{eq:phase10_support_q}
\frac{\partial c_i^{\mathrm{C},k}}{\partial q_j}=0
\quad \text{for } j\notin \mathcal S_q(i,k),
\end{equation}
\begin{equation}
\label{eq:phase10_support_T}
\frac{\partial c_i^{\mathrm{C},k}}{\partial T_j}=0
\quad \text{for } j\notin \mathcal S_T(i,k),
\end{equation}
and
\begin{equation}
\label{eq:phase10_support_card}
|\mathcal S_q(i,k)|\le C_q k,
\qquad
|\mathcal S_T(i,k)|\le C_T k
\end{equation}
for constants \(C_q,C_T>0\) independent of \(M\).
\end{assumption}

\begin{remark}
\label{rem:phase10_general_k_support_status}
A full constructive proof of Assumption~\ref{ass:phase10_general_k_support} for general \(k\) is \textbf{not completed at the present stage}.
This assumption is the general-\(k\) extension of the exact locality established in Phase~6 for the canonical raw Scheme~C map.
\end{remark}

Under these assumptions, we can derive the exact block-banded reverse pass.

Let
\[
J_{c\bar q}^{(k)}(\theta)
:=
\frac{\partial c^{\mathrm{C},k}}{\partial \bar q}(\theta)
\in \mathbb R^{2sM\times(M-1)},
\qquad
J_{cT}^{(k)}(\theta)
:=
\frac{\partial c^{\mathrm{C},k}}{\partial T}(\theta)
\in \mathbb R^{2sM\times M}.
\]

\begin{theorem}[Exact chain rule for the full framework]
\label{thm:phase10_chain_rule}
Assume Proposition~\ref{prop:phase10_K_C1} and Assumption~\ref{ass:phase10_general_k_C1}.
Then for
\[
J(\bar q,T):=K(c^{\mathrm{C},k}(\bar q,T),T),
\]
the gradients with respect to \(\bar q\) and \(T\) are
\begin{align}
\label{eq:phase10_grad_q}
\nabla_{\bar q}J
&=
\bigl(J_{c\bar q}^{(k)}\bigr)^\top g_c,\\
\label{eq:phase10_grad_T}
\nabla_T J
&=
\bigl(J_{cT}^{(k)}\bigr)^\top g_c
+
g_T.
\end{align}
\end{theorem}

\begin{proof}
This is the same chain-rule argument as in Phase~9.
Let \(dc\) be the differential of \(c^{\mathrm{C},k}(\bar q,T)\):
\[
dc
=
J_{c\bar q}^{(k)}\,d\bar q + J_{cT}^{(k)}\,dT.
\]
Since \(J(\bar q,T)=K(c(\bar q,T),T)\),
\[
dJ = g_c^\top dc + g_T^\top dT.
\]
Substituting the expression for \(dc\) yields
\[
dJ
=
g_c^\top J_{c\bar q}^{(k)}d\bar q
+
g_c^\top J_{cT}^{(k)}dT
+
g_T^\top dT.
\]
Collecting the coefficients of \(d\bar q\) and \(dT\) gives \eqref{eq:phase10_grad_q} and \eqref{eq:phase10_grad_T}.
\end{proof}

\begin{proposition}[Exact block-banded reverse accumulation]
\label{prop:phase10_banded_accum}
Under Assumption~\ref{ass:phase10_general_k_support}, the gradient components satisfy
\begin{equation}
\label{eq:phase10_banded_q}
\frac{\partial J}{\partial q_j}
=
\sum_{i:\, j\in \mathcal S_q(i,k)}
\left(\frac{\partial c_i^{\mathrm{C},k}}{\partial q_j}\right)^\top g_{c,i},
\end{equation}
and
\begin{equation}
\label{eq:phase10_banded_T}
\frac{\partial J}{\partial T_j}
=
g_{T,j}
+
\sum_{i:\, j\in \mathcal S_T(i,k)}
\left(\frac{\partial c_i^{\mathrm{C},k}}{\partial T_j}\right)^\top g_{c,i}.
\end{equation}
\end{proposition}

\begin{proof}
The dense formulas \eqref{eq:phase10_grad_q}--\eqref{eq:phase10_grad_T} expand blockwise as
\[
\frac{\partial J}{\partial q_j}
=
\sum_{i=1}^{M}
\left(\frac{\partial c_i^{\mathrm{C},k}}{\partial q_j}\right)^\top g_{c,i},
\qquad
\frac{\partial J}{\partial T_j}
=
g_{T,j}
+
\sum_{i=1}^{M}
\left(\frac{\partial c_i^{\mathrm{C},k}}{\partial T_j}\right)^\top g_{c,i}.
\]
By Assumption~\ref{ass:phase10_general_k_support}, all terms outside the support sets vanish exactly.
Therefore the sums reduce to \eqref{eq:phase10_banded_q} and \eqref{eq:phase10_banded_T}.
\end{proof}

\begin{corollary}[Dense checker and sparse implementation are mathematically identical]
\label{cor:phase10_dense_sparse_identical}
Under Assumption~\ref{ass:phase10_general_k_support}, a dense reverse pass based on \eqref{eq:phase10_grad_q}--\eqref{eq:phase10_grad_T} and a sparse reverse pass based on \eqref{eq:phase10_banded_q}--\eqref{eq:phase10_banded_T} produce exactly the same gradient.
\end{corollary}

\begin{proof}
The sparse pass removes only terms that are identically zero by \eqref{eq:phase10_support_q}--\eqref{eq:phase10_support_T}. Hence the sums are equal term by term.
\end{proof}

\subsection{Complexity of full objective and gradient evaluation}
\label{subsec:phase10_complexity}

We now give the full framework complexity theorem under an explicit computational model.

\begin{assumption}[Per-segment assembly cost]
\label{ass:phase10_assembly_cost}
There exists a constant \(C_{\mathrm{asm}}>0\) such that for each segment \(i\):
\begin{enumerate}
    \item the coefficient block \(c_i^{\mathrm{C},k}\) can be assembled in \(O(k^2)\) time;
    \item all nonzero Jacobian blocks
    \[
    \left\{
    \frac{\partial c_i^{\mathrm{C},k}}{\partial q_j}
    : j\in \mathcal S_q(i,k)
    \right\},
    \qquad
    \left\{
    \frac{\partial c_i^{\mathrm{C},k}}{\partial T_j}
    : j\in \mathcal S_T(i,k)
    \right\}
    \]
    can be assembled in total \(O(k^2)\) time;
    \item each sampled penalty evaluation uses \(O(1)\) sample points per segment and each sample contribution is \(O(1)\).
\end{enumerate}
\end{assumption}

\begin{remark}
\label{rem:phase10_assembly_cost_status}
A low-level proof that the chosen implementation achieves Assumption~\ref{ass:phase10_assembly_cost} for every general-\(k\) raw Scheme~C solver is \textbf{not completed at the present stage}.
This assumption is the mathematically clean interface between the BLOM local solver and the framework-level complexity theorem.
\end{remark}

\begin{theorem}[Framework-level complexity theorem]
\label{thm:phase10_complexity}
Under Assumptions~\ref{ass:phase10_general_k_support} and \ref{ass:phase10_assembly_cost}, one evaluation of:
\begin{enumerate}
    \item the coefficients \(c^{\mathrm{C},k}(\bar q,T)\),
    \item the full objective \(K(c,T)\),
    \item the block gradients \(g_c(c,T)\) and \(g_T(c,T)\),
    \item and the parameter gradient \(\nabla_\theta J\),
\end{enumerate}
can be performed in
\begin{equation}
\label{eq:phase10_time_complexity}
O(k^2 M)
\end{equation}
time and
\begin{equation}
\label{eq:phase10_memory_complexity}
O(kM)
\end{equation}
memory.
\end{theorem}

\begin{proof}
We decompose the computation into four steps.

\paragraph{Step 1: coefficient and local Jacobian assembly.}
By Assumption~\ref{ass:phase10_assembly_cost}, each segment requires \(O(k^2)\) time to assemble its coefficient block and all its nonzero Jacobian blocks.
Over all \(M\) segments, this costs \(O(k^2M)\).

\paragraph{Step 2: objective and block-gradient evaluation.}
The control term requires \(O(1)\) work per segment because the block size \(2s\) is fixed.
The time regularizer and time penalty are also \(O(1)\) per segment.
Obstacle and dynamics penalties use \(O(1)\) sample points per segment by assumption, so their total cost is \(O(M)\).
Boundary terms are \(O(1)\).
Hence the full objective and all block gradients \(g_c,g_T\) are computed in \(O(M)\), which is dominated by \(O(k^2M)\).

\paragraph{Step 3: sparse backward accumulation.}
For each coefficient block \(c_i\), only \(O(k)\) parameter blocks interact with it by Assumption~\ref{ass:phase10_general_k_support}.
Since each block-vector product involves fixed block size \(2s\), the work per segment is \(O(k)\).
Thus the reverse accumulation cost is \(O(kM)\), which is again dominated by \(O(k^2M)\).

\paragraph{Step 4: memory.}
For each segment, the number of stored nonzero Jacobian blocks is \(O(k)\), and each block has fixed size \(2s\times 1\).
Therefore the total storage for all sparse Jacobian blocks is \(O(kM)\).
The coefficient blocks and gradient blocks are \(O(M)\), hence the overall memory cost remains \(O(kM)\).

Combining these steps proves \eqref{eq:phase10_time_complexity} and \eqref{eq:phase10_memory_complexity}.
\end{proof}

\begin{remark}
\label{rem:phase10_complexity_scope}
Theorem~\ref{thm:phase10_complexity} is the correct framework-level complexity statement for Phase~10.
It does not yet prove the low-level numerical linear algebra complexity of every possible raw Scheme~C solver; that implementation-specific layer remains part of future work.
\end{remark}

\subsection{Optimization algorithm on the unconstrained variable \texorpdfstring{$\xi=(\bar q,\tau)$}{xi}}
\label{subsec:phase10_optimizer}

With the time reparameterization in place, the unconstrained objective is
\[
\widetilde J(\bar q,\tau)
=
J(\bar q,T(\tau)).
\]

\begin{proposition}[Exact gradient in the reparameterized variable]
\label{prop:phase10_tau_gradient}
Let \(\widetilde J(\bar q,\tau)=J(\bar q,T(\tau))\).
Then
\begin{equation}
\label{eq:phase10_tau_grad_q}
\nabla_{\bar q}\widetilde J
=
\nabla_{\bar q}J,
\end{equation}
and
\begin{equation}
\label{eq:phase10_tau_grad_tau}
\nabla_{\tau}\widetilde J
=
D_{\sigma}(\tau)\,\nabla_T J,
\qquad
D_{\sigma}(\tau):=\operatorname{diag}(\sigma(\tau_1),\dots,\sigma(\tau_M)).
\end{equation}
\end{proposition}

\begin{proof}
The first identity is immediate because \(\bar q\) is unchanged by the reparameterization.
For the second identity, apply the chain rule componentwise:
\[
\frac{\partial \widetilde J}{\partial \tau_i}
=
\frac{\partial J}{\partial T_i}\frac{dT_i}{d\tau_i}.
\]
By Proposition~\ref{prop:phase10_softplus},
\[
\frac{dT_i}{d\tau_i}=\sigma(\tau_i).
\]
Therefore
\[
\frac{\partial \widetilde J}{\partial \tau_i}
=
\sigma(\tau_i)\frac{\partial J}{\partial T_i},
\]
which is exactly \eqref{eq:phase10_tau_grad_tau}.
\end{proof}

We now consider a standard backtracking line-search step on \(\xi=(\bar q,\tau)\).

\begin{theorem}[Local descent with Armijo backtracking]
\label{thm:phase10_armijo}
Assume \(\widetilde J\in C^1\) and that \(\nabla \widetilde J\) is Lipschitz continuous on a convex neighborhood \(U\subset\mathbb R^{M-1+M}\).
Let
\[
d(\xi):=-\nabla \widetilde J(\xi).
\]
Fix constants
\[
0<\beta<1,
\qquad
0<c_{\mathrm{A}}<1.
\]
Then for every \(\xi\in U\) with \(\nabla \widetilde J(\xi)\neq 0\), there exists an integer \(m\ge 0\) such that the step size
\[
\alpha=\beta^m
\]
satisfies the Armijo inequality
\begin{equation}
\label{eq:phase10_armijo}
\widetilde J(\xi+\alpha d(\xi))
\le
\widetilde J(\xi)-c_{\mathrm{A}}\alpha\|\nabla \widetilde J(\xi)\|_2^2.
\end{equation}
\end{theorem}

\begin{proof}
Because \(\nabla\widetilde J\) is Lipschitz, the standard descent lemma implies that for sufficiently small \(\alpha>0\),
\[
\widetilde J(\xi-\alpha \nabla \widetilde J(\xi))
\le
\widetilde J(\xi)
-
\alpha\left(1-\frac{L\alpha}{2}\right)\|\nabla \widetilde J(\xi)\|_2^2,
\]
where \(L\) is a Lipschitz constant of \(\nabla\widetilde J\) on \(U\).
Choose \(\alpha>0\) small enough so that
\[
1-\frac{L\alpha}{2}\ge c_{\mathrm{A}}.
\]
Then
\[
\widetilde J(\xi-\alpha \nabla \widetilde J(\xi))
\le
\widetilde J(\xi)-c_{\mathrm{A}}\alpha\|\nabla \widetilde J(\xi)\|_2^2.
\]
Since the geometric sequence \(\beta^m\) tends to zero as \(m\to\infty\), there exists \(m\) such that \(\alpha=\beta^m\) is small enough to satisfy the above inequality.
\end{proof}

\begin{remark}[Marked limitation]
\label{rem:phase10_global_conv_status}
A global convergence theorem for the full nonconvex optimizer, including obstacle and dynamics penalties, is \textbf{not completed at the present stage}.
Phase~10 proves the correct local descent statement under standard smoothness assumptions, which is the mathematically appropriate guarantee for a framework chapter.
\end{remark}

\subsection{Benchmark protocol and baseline comparisons}
\label{subsec:phase10_benchmark}

The purpose of the benchmark layer is to compare the raw Scheme~C framework against representative baselines without changing the underlying task definition.

\paragraph{Baselines.}
The primary baselines are:
\begin{enumerate}
    \item the global MINCO reference baseline;
    \item the shared-state Scheme~A baseline;
    \item a heuristic local interpolation baseline.
\end{enumerate}

\paragraph{Evaluation metrics.}
For each method, define:
\begin{equation}
\label{eq:phase10_metric_obj}
\Delta_{\mathrm{obj}}
:=
\frac{K_{\mathrm{final}}-K_{\mathrm{ref}}}{|K_{\mathrm{ref}}|+\varepsilon},
\end{equation}
\begin{equation}
\label{eq:phase10_metric_time}
\Delta_{\mathrm{time}}
:=
\text{wall-clock time per iteration},
\end{equation}
\begin{equation}
\label{eq:phase10_metric_mem}
\Delta_{\mathrm{mem}}
:=
\text{peak memory},
\end{equation}
and, if relevant,
\begin{equation}
\label{eq:phase10_metric_violation}
\Delta_{\mathrm{viol}}
:=
\text{sampled constraint violation}.
\end{equation}

\paragraph{Sweeps.}
The benchmark suite should vary at least:
\begin{enumerate}
    \item the number of segments \(M\),
    \item the bandwidth \(k\),
    \item obstacle density or penalty strength,
    \item time-range spread,
    \item initialization quality.
\end{enumerate}

\paragraph{Ablations.}
A complete ablation suite should include:
\begin{enumerate}
    \item \(k\)-sweep ablations;
    \item term-ablation studies removing \(K_{\mathrm{obs}}\), \(K_{\mathrm{dyn}}\), or \(K_{\mathrm{time}}\);
    \item raw Scheme~C versus baseline comparisons under matched task settings.
\end{enumerate}

\begin{remark}
\label{rem:phase10_benchmark_scope}
The benchmark protocol is part of the framework definition, not a theorem.
Its role is to make the framework experimentally falsifiable and directly comparable to established baselines.
\end{remark}

\subsection{What is proved and what remains open in Phase~10}
\label{subsec:phase10_status}

We summarize the exact status of this phase.

\paragraph{Fully proved in this chapter.}
\begin{enumerate}
    \item the smooth positivity-preserving time reparameterization and its derivative;
    \item the \(C^1\) well-posedness of the full objective layer \(K(c,T)\);
    \item the exact block formulas for \(\partial K/\partial c\) and \(\partial K/\partial T\);
    \item the exact chain-rule gradient formulas for \(J(\bar q,T)\);
    \item the exact sparse reverse accumulation formulas under explicit support assumptions;
    \item the framework-level complexity theorem \(O(k^2M)\) time and \(O(kM)\) memory under explicit assembly-cost assumptions;
    \item the exact gradient formula in the unconstrained variable \(\xi=(\bar q,\tau)\);
    \item the existence of an Armijo descent step under standard smoothness assumptions.
\end{enumerate}

\paragraph{Not completed at the present stage.}
\begin{enumerate}
    \item a full proof of general-\(k\) differentiability and exact local support of the raw Scheme~C map;
    \item a constructive proof that the chosen implementation always realizes the assumed \(O(k^2)\) per-segment assembly cost;
    \item a global convergence theorem for the full nonconvex optimization process;
    \item feasibility / cost-theoretic bridges back to the exact global MINCO mother problem.
\end{enumerate}

Accordingly, the mathematically correct summary of Phase~10 is:
\[
\boxed{
\text{raw Scheme~C has now been elevated from a differentiable representation to a complete block-banded space-time optimization framework,}
}
\]
but
\[
\boxed{
\text{the strongest global optimization and feasibility-theoretic claims remain outside the scope of the present phase.}
}